%!TEX root = main.tex

\smallskip\noindent
{\bf Benchmarks}. To extensively evaluate the OCaml FIRRTL compiler,
we mainly collect FIRRTL circuits from the llvm/circt \footnote{\url{https://github.com/llvm/circt/tree/main}.} and the Chisel Book--Digital Design with Chisel~\cite{chisel:book}.

There are 121 FIRRTL circuits inside the ``llvm/circt/test/firtool'' directory
of circt, among which 77 FIRRTL circuits are excluded because they are not self-contained (i.e., dependent upon external modules),
or contain undefined semantics for memories, or have no output ports, or have illegal connections, or cannot be handled by our compiler due to probe and instance statements.
Thus, we consider the remaining 44 FIRRTL circuits from circt.

There are 51 Chisel examples in Scala inside the ``schoeberl/chisel-book/src/main/scala'' directory of the Chisel Book.
These 51 Chisel examples are transformed into 62 FIRRTL circuits by applying the convertor \textsf{ChiselStage} within the Chisel compiler framework
to each subclass of the class $\tt Module$ in Scala. (Note that some Chisel examples contain multiple such subclasses, thus yield multiple FIRRTL circuits.)

%For those from Treadle are 55 LoFIRRTL circuits inside the ``chipsalliance/treadle/samples'' directory.
%We consider those LoFIRRTL circuits because Treadle is more likely to execute correctly on them so that
%the correctness testing of the formalized operational semantics has high confidence.
%Among 55 LoFIRRTL circuits, we exclude 18 LoFIRRTL circuits because they violate the syntax of
%LoFIRRTL (i.e., dependency of external modules), undefined semantics for memories, have no output ports or
%cannot be handled by Treadle due to cyclical connection of ports.
%Thus, we collected 37 LoFIRRTL circuits from Treadle, among them 3 circuits have instof statement,  which means a nested module structure.
%Beside, we generated 76 LoFIRRTL circuits from book project of Digital Design with Chisel, which containing 51 Chisel examples inside the ``schoeberl/chisel-book/src/main/scala'' directory. The transformation is executed by Chisel stage for each class extends from class $\tt Module$, as it is legal to define multi-class in one Scala program, it is possible to execute 76 LoFIRRTL circuits from 51 Chisel programs. The generated LoFIRRTL circuits including 66 single module circuits and 10 with instance structure.
Besides the above FIRRTL circuits, we also collect some issues reported as bug of these three passes on the github and designed 7 examples to test the correctness of Coq compiler.

In summary, we use 113 (44+62+7) FIRRTL circuits all of which can be handled by firtool.

\smallskip\noindent
{\bf Testing method}.
Since we want to compare the equivalence of circuit structure generated by our FIRRTL OCaml compiler and firtool, we first produce LoFIRRTL program of MLIR syntax and use firtool to transform it verilog. Then we use yosys \footnote{\url{https://github.com/YosysHQ/yosys}.} to transform verilog programs to aiger files and compare the equivalence of our compile results and the official ones by berkeley-abc \footnote{\url{https://github.com/berkeley-abc/abc}.}.

\begin{table}[!t]
\caption{Results of conformance testing}\small %\footnotesize\vspace*{-3mm}
\label{tab:conformancetesting}
\scalebox{0.95}{\begin{tabular}{|c|c|c|c|c|c|c|}
  \hline
\multirow{2}{*}{\bf Benchmark suit} & \multirow{2}{*}{\bf \#LOC} & \multicolumn{2}{c|}{\bf Compilation time (ms)}&\multirow{2}{*}{\bf Speedup}  & \multirow{2}{*}{\bf Verification time} & \multirow{2}{*}{\bf Result} \\ \cline{3-4}
                              &                        & {\bf firtool}  &  {\bf Ours}  &     &  \\ \hline \hline
      circt (44 circuits) & 36.43  & 1.10  & 0.32  & 2.50  & 0.00 & \ding{51} \\ \hline
         Chisel Book (62 circuits) & 79.13 & 9.89  & 7.68  &  0.29 & 0.00 & \ding{51} \\ \hline
        issue (7 circuits) & 44.29 & 1.35  & 0.044  & 29.68 & 0.00 & \ding{55}(1) \\ \hline \hline
        \textbf{Average} &  \textbf{60.35} & 	\textbf{5.94} & 	\textbf{4.34} & 	\textbf{0.37} & \textbf{0.00} &  \\ \hline
 \end{tabular}}
\end{table}
 	 	 	

\smallskip\noindent
{\bf Testing results}.
The summary of conformance testing is reported in Table~\ref{tab:conformancetesting}.
Column (\#LOC) shows the average number of locations for each benchmark suit.
Column (firtool) shows the average compilation time $t_1$ of firtool in millisecond (ms) for each benchmark suit.
Column (Ours) shows the average compilation time $t_2$ of our FIRRTL OCaml compiler in ms for each benchmark suit.
Column (Speedup) is calculated via $\frac{t_1-t_2}{t_2}$ for each benchmark suit.
Column (Verification time) shows the average verification time $t_v$ of equivalence checking by berkeley-abc in millisecond (ms) for each benchmark suit.
Column (Result) shows the comparison result of circt structure between firtool and our FIRRTL OCaml compiler under berkeley-abc,
where \ding{51} indicates that the circuits produced by two tools have the same structure and \ding{55}(n) indicates
that the circuits of $n$ benchmarks produced by two tools are different.
The last row (Average) shows the average results over all 113 benchmarks.
Detailed results are given in the supplementary material due to space limits.

Overall, our FIRRTL OCaml compiler successfully executes all the 106 benchmarks from circt and chiselbook and produces the same circuit structure as firtool. In the official circt database, there are five bugs associated with these three passes, four of which is about inferwidths and one is about expandconnects. Four of these have been resolved, ensuring they can be executed correctly in the latest firtool. One of the bugs related to inferwidths remains unresolved. This issue arises from the fact that firtool uses constraint solver to infer bit widths, and their algorithm detected a loop in the circuit. However, the circuit does not form a loop and should be solvable. Our Coq compiler implement a verified topological sorting algorithm so that LoFIRRTL can be successfully generated.

In terms of efficiency, our FIRRTL OCaml compiler is faster (0.37 times on average) than firtool. Only 3 benchmarks from ChiselBook cost more compilation time than firtool. 
After an in-depth analysis, we found that in \texttt{DeepExecution} with 1220 lines our compiler has 3 times compilation time than firtool. Inferwidths pass costs 30 times the compilation time than firtool, because the time cost in topological sorting algorithm increase as the dependency graph involve more components. We plan to implement a constraint solver as well and provide these two methods as an option.

